\begin{abstract}
\thispagestyle{empty}

Выпускная квалификационная работа (далее ВКР) на тему: <<Разработка программного комплекса на основе мультиагентной архитектуры с~гибридной моделью интеграции для управления цифровым присутствием субъектов малого и среднего бизнеса>>.

\textbf{Объект исследования}~--- процесс управления цифровым присутствием субъектов малого и среднего бизнеса (МСБ) на внешних платформах.

\textbf{Предмет исследования}~--- автоматизация управления цифровым присутствием на основе мультиагентной архитектуры с~гибридной моделью интеграции (программные интерфейсы приложений, API~+ роботизированная автоматизация процессов, RPA).

\textbf{Цель работы}~--- разработка и~реализация программного комплекса на основе мультиагентной архитектуры с~гибридной моделью интеграции, обеспечивающего автоматизацию управления цифровым присутствием субъектов МСБ на нескольких внешних площадках.

\textbf{Методы исследования:} сравнительный анализ существующих решений методом взвешенных критериев, объектно-ориентированное проектирование, моделирование бизнес-процессов (AS-IS / TO-BE), инфологическое и даталогическое моделирование данных, микросервисная декомпозиция, функциональное тестирование по критическому пути пользователя.

\textbf{Результаты работы.} Проведён анализ предметной области и~семи существующих решений трёх классов (управление взаимоотношениями с~клиентами, планирование публикаций, мониторинг репутации). Сформулированы функциональные и~нефункциональные требования к~системе. Спроектирована и~реализована микросервисная система из шести компонентов и~разделяемой библиотеки. Реализованы три платформенных агента (Telegram, VK, Яндекс.Бизнес), из которых два работают через официальные API платформ, а~один~--- через браузерную автоматизацию (Playwright). Разработан пользовательский интерфейс в~виде одностраничного веб-приложения на базе Next.js~14 с~10~экранными формами.

\textbf{Объём работы:} 61~страница, 19~рисунков, 18~таблиц, 31~источник.

\textbf{Ключевые слова:} мультиагентная система, цифровое присутствие, малый и~средний бизнес, оркестрация, браузерная автоматизация, роботизированная автоматизация процессов (RPA), микросервисная архитектура, большая языковая модель (LLM), Go, Next.js.

\end{abstract}
